% Miscellaneous - R(sigma,delta) with two crossings of lambda^R.
% GENERATED BY generate_figures.py - edit that file, not this one.
% delta*_1=0.9098 (dR/ddelta=+0.1030), delta*_2=1.8496 (dR/ddelta=-0.1181)
% R runs over [0.1045, 0.2306]; two crossings for any lambda^R in (0.1654, 0.2306)
\pgfplotsset{miscacc/.style={
  width=12.0cm, height=7.4cm,
  xmin=-0.10, xmax=2.57, ymin=0.075, ymax=0.266,
  axis x line=bottom, axis y line=left,
  axis line style={-{Stealth[length=6pt]},black}, clip=false,
  xlabel={$\delta$}, ylabel={$\mathcal{R}(\sigma,\delta)$},
  every axis x label/.style={at={(current axis.south east)},anchor=west,xshift=9pt},
  every axis y label/.style={at={(current axis.north west)},anchor=south,yshift=7pt},
  tick label style={font=\footnotesize}, label style={font=\footnotesize},
  title style={font=\footnotesize, yshift=2pt},
  tick style={black},
  xtick={0.9098,1.8496}, xticklabels={$\delta^{*}_{1}$,$\delta^{*}_{2}$},
  ytick={0.2000}, yticklabels={$\lambda^{R}$},
}}

\begin{figure}[htbp]
\centering
\begin{tikzpicture}
\begin{axis}[miscacc]
  \addplot[black,dashed,thick,samples=2,domain=-0.10:2.57] {0.2000};
  \draw[gray,dashed] (axis cs:0.90984,0.075) -- (axis cs:0.90984,0.20000);
  \draw[gray,dashed] (axis cs:1.84963,0.075) -- (axis cs:1.84963,0.20000);
  \addplot[blue,very thick,smooth] coordinates {
(0.0500,0.1045) (0.0807,0.1074) (0.1114,0.1105) (0.1421,0.1135) (0.1728,0.1167)
(0.2035,0.1199) (0.2342,0.1232) (0.2649,0.1266) (0.2956,0.1300) (0.3263,0.1334)
(0.3569,0.1369) (0.3876,0.1404) (0.4183,0.1439) (0.4490,0.1475) (0.4797,0.1511)
(0.5104,0.1547) (0.5411,0.1583) (0.5718,0.1619) (0.6025,0.1655) (0.6332,0.1691)
(0.6639,0.1727) (0.6946,0.1762) (0.7253,0.1798) (0.7560,0.1833) (0.7867,0.1867)
(0.8174,0.1901) (0.8481,0.1935) (0.8788,0.1968) (0.9095,0.2000) (0.9402,0.2031)
(0.9708,0.2061) (1.0015,0.2090) (1.0322,0.2118) (1.0629,0.2144) (1.0936,0.2169)
(1.1243,0.2193) (1.1550,0.2214) (1.1857,0.2234) (1.2164,0.2252) (1.2471,0.2267)
(1.2778,0.2280) (1.3085,0.2291) (1.3392,0.2299) (1.3699,0.2304) (1.4006,0.2306)
(1.4313,0.2305) (1.4620,0.2301) (1.4927,0.2293) (1.5234,0.2283) (1.5541,0.2269)
(1.5847,0.2252) (1.6154,0.2232) (1.6461,0.2210) (1.6768,0.2184) (1.7075,0.2155)
(1.7382,0.2125) (1.7689,0.2092) (1.7996,0.2058) (1.8303,0.2023) (1.8610,0.1987)
(1.8917,0.1950) (1.9224,0.1913) (1.9531,0.1877) (1.9838,0.1842) (2.0145,0.1807)
(2.0452,0.1774) (2.0759,0.1742) (2.1066,0.1711) (2.1373,0.1682) (2.1680,0.1654)};
\end{axis}
\end{tikzpicture}
\caption{The regular acceptance rate $\mathcal{R}(\sigma,\delta)$ of Vertical\_Model\_4 against the normalised opportunity cost $\delta=(\hat{\theta}^T-\hat{\theta}^R)/\sigma$. It is hump shaped, so the payline $\lambda^R$ is crossed twice: $\delta^*_1$ is stable ($\partial\mathcal{R}/\partial\delta>0$) and $\delta^*_2$ is not.}
\label{fig:misc_acc}
\end{figure}
